\subsection{膜结构}\label{sec:13c}
现在假设 $M$ 带有二次复体积形式；在这一维数中，它就是处处非零的 $C^\infty$ 二次微分 $\eta_M^2:TM\otimes_{\CC} TM\to\CC$。相应平方相位映射可视为实射影化切丛的平凡化 $\alpha_M:\RP(TM)=\operatorname{Gr}(TM)\to S^1$。一种好的直观图像来自实子丛 $\xi_M=\alpha_M^{-1}(1)$，它给出曲面上的无定向叶状结构。用这一语言，$M$ 中的 Lagrangian 膜就是嵌入闭曲线 $L$，配有
\begin{itemize}
\item 拉回丛 $\pi^*(TM|L)$ 的实子丛 $H$，其中 $\pi:[0;1]\times L\to L$ 是投影，满足 $H|\{0\}\times L=\xi_M|L$、$H|\{1\}\times L=TL$；
\item 实线丛 $\beta\to L$。
\end{itemize}
$H$ 是 $TM|L$ 的子丛同伦；从几何看，就是同时旋转 $L$ 的所有切空间，使之与 $\xi_M$ 平行的一种方式。这种同伦存在的障碍是 $L$ 相对于 $\xi_M$ 的整体旋转数；要求该数为零，再次排除了零同调的简单闭曲线。为与一般分次概念联系，注意：给定这种 $H$ 后，存在唯一函数 $a:[0;1]\times L\to\R$，满足 $a(0,\allowbreak x)=0$ 与 $\exp(2\pi ia(r,\allowbreak x))=\alpha_M(H_{r,\allowbreak x})$，于是 $\alpha^\#(x)=a(1,\allowbreak x)$ 是 $L$ 的一个分次。至于 $\beta$，由\stref{例题}{11.20}已知，它等价于 Pin 结构的一个选择。\znote{195}{曲面上的膜结构可以完全用初等数据描述：$H$ 是从 $\xi_M|L$ 到 $TL$ 的子丛同伦，旋转角的累计给出分次（见 (13.6) 前的构造）；$\beta$ 则代替 Pin 结构，因为在实一维情形，$TL$ 上的 Pin 结构恰好等价于一个实线丛（例题 11.20）。这解释了 13c 为何能完全避开 12a 的 $\operatorname{Gr}(TM)^\#$ 语言：二维时 $\alpha_M:\RP(TM)\to S^1$ 已是射影化切丛的平凡化，$\xi_M$ 就是它的水平叶状结构。由 $H$ 与 $\beta$ 还能看出，膜的自同构群恰是保持这两项数据的复形，这与 12i 中的 $\operatorname{Aut}^\#$ 相呼应。[来源：Seidel, Fukaya categories and Picard–Lefschetz theory, §13c 与例题 11.20。]}

令 $\symbfit{L}^\#$ 为处于一般位置的一族正合 Lagrangian 膜 $L_k^\#=(L_k,\allowbreak H_k,\allowbreak \beta_k)$，$1\leq k\leq m$。现在为先前构造加入分次与符号，得到组合定义的有向 $\Ai$ 范畴。首先，为每个交点 $y\in L_{i_0}\cap L_{i_1}$（$i_0<i_1$）赋予绝对指标 $i(y)$ 和定向空间 $o(y)$。将同伦 $H_{i_0}(1-\cdot,\allowbreak y)$ 与 $H_{i_1}(\cdot,\allowbreak y)$ 组合为 $\RP(TM_y)$ 中从 $(TL_{i_0})_y$ 到 $(TL_{i_1})_y$ 的路径。令 $\measuredangle(y)\in\R\setminus\pi\Z$ 为该路径旋转的总角度，并设
\begin{equation}\label{eq:13.6}
\begin{aligned}
i(y)&=\lfloor\measuredangle(y)/\pi\rfloor+1,\\
o(y)&=\Hom(\beta_{i_0,y},\beta_{i_1,y})\otimes(TL_{i_1})_y^{\otimes i(y)}.
\end{aligned}
\end{equation}
然后如\eref{12.16}定义态射空间 $CF^*(L_{i_0}^\#,\allowbreak L_{i_1}^\#)$。接着，取角点位于 $y_0,\allowbreak \ldots,\allowbreak y_d$ 的浸入多边形 $\tilde u:\Delta^{d+1}\to M$。用\eref{13.6}作简单计算可见，这种多边形的存在蕴含等式 $i(y_0)=i(y_1)+\cdots+i(y_d)+2-d$。用 $\Delta^{d+1}$ 在各角点的边界定向，赋予 $(TL_{i_d})_{y_0}$ 和 $(TL_{i_k})_{y_k}$、$k=1,\allowbreak \ldots,\allowbreak d$ 定向，如\sr{fig:13.1}{图~13.1}中的箭头所示；注意 $y_0$ 处与 $y_1,\allowbreak \ldots,\allowbreak y_d$ 处的约定不同。沿 $\Delta^{d+1}$ 各边的平行移动定义同构 $\beta_{i_k,\allowbreak y_k}\to\beta_{i_k,\allowbreak y_{k+1}}$ 及 $\beta_{i_d,\allowbreak y_d}\to\beta_{i_d,\allowbreak y_0}$。由\eref{13.6}可见，这两个初等事实
\SourcePage{204}
结合给出同构 $o(y_d)\otimes\cdots\otimes o(y_1)\to o(y_0)$，记为 $\tilde c_{\tilde u}$。将全部规范化 $|\tilde c_{\tilde u}|_{\K}$ 相加，定义
\[
\tilde\mu^d:CF^*(L_{i_{d-1}}^\#,L_{i_d}^\#)\otimes\cdots\otimes CF^*(L_{i_0}^\#,L_{i_1}^\#)
\longrightarrow CF^{*+2-d}(L_{i_0}^\#,L_{i_d}^\#).
\]
一个基于具有一个凹角的多边形分解的初等论证表明，这些映射满足 $\Ai$ 结合方程\eref{1.2}。\znote{196}{这里的论证是“带一个凹角的多边形分解”：沿一条弦把多边形切成两块，正好对应把一个余维一边界面拆成两个低阶模空间，凹角保证切开后各部分仍是凸多边形。于是 $d$ 阶关系归结为 $d-1$ 阶与更低阶关系的组合，与 12g 中 Deligne–Mumford 紧化的论证完全平行，只是把“模空间边界”换成了“多边形剖分”。这类只数多边形的范畴是纯组合的，文献中正是用它在带二次微分的曲面上研究稳定性条件。[来源：arXiv:1409.8611，Haiden–Katzarkov–Kontsevich, Flat surfaces and stability structures，摘要（2026-11-15 访问）；并见原书 §13c。]}将由此定义的有向 $\Ai$ 范畴记为 $\widetilde{\F}(\symbfit{L}^\#)^\to$。
\begin{example}\label{st:13.5}
考虑 $M=S^1\times[-1;1]$；\sr{fig:13.2}{图~13.2}把它画成两个竖直边相粘合的正方形，并画出其上的曲线 $L_0,\allowbreak L_1$。使用对应于水平直线叶状结构 $\xi_M$ 的常值复体积形式。选择从 $\xi_M|L_0$ 到 $TL_0$ 的平凡同伦。另一条曲线中，$\xi_M|L_1$ 与 $TL_1$ 的夹角处处小于 $\pi/2$，所以二者之间有唯一的“短”同伦，采用这一同伦。最后取平凡线丛 $\beta_0,\allowbreak \beta_1$。则 $i(y_0)=0$，
\SourceFigure{figure-13-2}{0.35099962}{13.2}
$i(y_1)=1$，$o(y_0)=\R$，且此识别是典范的，因为指标为偶数；另有 $o(y_1)=(TL_1)_{y_1}$。\sr{fig:13.1}{图~13.1}的约定说明，计算二边形 $u$ 对边界算子的贡献时，应使用\sr{fig:13.2}{图~13.2}所示的定向识别 $o(y_1)\cong\R$。另一个二边形 $v$ 要求相反识别，所以二者贡献抵消，得到
\[
HF^0(L_0^\#,L_1^\#)=\K,\qquad HF^1(L_0^\#,L_1^\#)=|(TL_1)_{y_1}|_{\K}.
\]
因为 $L_0,\allowbreak L_1$ 作为正合 Lagrangian 膜同痕，一般理论说明存在典范同构 $HF^*(L_0^\#,\allowbreak L_1^\#)\cong H^*(L_1;\K)$。确实，$H^0(L_1;\K)\cong\K$ 是典范的，而 $H^1(L_1;\K)$ 的自然生成元恰是 $L_1$ 的定向，因此 $H^1(L_1;\K)\cong|(TL_1)_{y_1}|_{\K}$。鼓励读者试验改变膜结构的效果，并将结果与\eref{12.12}、\eref{12.14}比较。
\end{example}

下面将上述组合构造与真正 Fukaya 范畴的有向子范畴 $\F(\symbfit{L}^\#)^\to\subset\F(M)$ 联系。对象显然
\SourcePage{205}
相同，态射空间也相同，因为\eref{13.6}只是\eref{11.35}的改写；如先前所见，零维模空间的点 $(r,\allowbreak u)\in\M^{d+1}(y_0,\allowbreak \ldots,\allowbreak y_d)$ 与浸入 $(d+1)$ 边形 $\tilde u$ 的重参数化等价类之间存在典范双射。唯一差别是：按\secref{12g}的一般方法构造 $\Ai$ 结构映射 $\mu^d$ 时，使用的是通过行列式线粘合论证在分析上定义的同构 $c_{r,\allowbreak u}$，见\eref{12.23}。更确切地，$\mu^d$ 由全部贡献 $|(-1)^\dagger c_{r,\allowbreak u}|_{\K}$ 相加得到，其中 $\dagger$ 如\eref{12.24}。据此，以 $\gamma^d(\tilde u)\in\{\pm1\}$ 表示 $\tilde c_{\tilde u}$ 与 $(-1)^\dagger c_{r,\allowbreak u}$ 之间的符号差。

断言 $\gamma^d(\tilde u)=\gamma^d(k_d,\allowbreak \ldots,\allowbreak k_1)$ 实际只由相关交点的绝对指标 $k_j=i(y_j)$ 决定。首先，符号差只依赖浸入多边形附近几何数据的局部行为；其次，若将 $M$ 换成某个覆盖，取 $L_{i_k}$ 的原像并任取多边形的一个提升，符号差不变。由此可将符号的确定归结到多边形嵌入而非仅浸入的情形。此时可沿各边平凡化线丛 $\beta_{i_k}$，说明 $\gamma^d(\tilde u)$ 对所有线丛选择都相同。剩下的局部拓扑数据只有绝对指标，故断言成立。\znote{197}{$\gamma^d$ 只依赖指标，理由有三层：其一，符号差是局部量，只取决于多边形附近的几何数据；其二，覆盖不改变符号差，故可把浸入提升为嵌入；其三，嵌入后可沿各边平凡化诸 $\beta_{i_k}$，线丛的选取不再影响符号。剩下的只有各角点的绝对指标 $k_j=i(y_j)$。这与 12g 中 $\dagger$、$\ddagger$、$\diamond$ 只依赖指标的现象同源：定向空间带来的符号最终都由 Maslov 指标决定。[来源：Seidel, Fukaya categories and Picard–Lefschetz theory, §13c；并与 §12g 的符号计算对照。]}

还有一项先验信息：$\mu^d$ 与 $\tilde\mu^d$ 均满足 $\Ai$ 结合方程。这对 $\gamma^d$ 施加相当强的限制。例如，考察\sr{fig:13.3}{图~13.3}(i) 所示的分解，可见
\[
\begin{aligned}
&\gamma^d(k_d,\ldots,k_{n+1},k_n+1,\ldots,k_1)\gamma^1(k_n)\\
&\qquad=\gamma^d(k_d,\ldots,k_{n+1}+1,k_n,\ldots,k_1)\gamma^1(k_{n+1}),\\
&\gamma^d(k_d+1,k_{d-1},\ldots,k_1)\gamma^1(k_d)\\
&\qquad=\gamma^d(k_d,\ldots,k_1)\gamma^1(k_1+\cdots+k_d+2-d).
\end{aligned}
\]
\SourceFigure{figure-13-3}{0.4880764}{13.3}
作初等变换，得到
\[
\frac{\gamma^d(k_d,\ldots,k_1)}{\gamma^d(0,\ldots,0)}
=\frac{\Gamma(k_1)\Gamma(k_2)\cdots\Gamma(k_d)\Gamma(k_1+\cdots+k_d+2-d)}{\Gamma(2-d)},
\]
\SourcePage{206}
其中 $k\geq0$ 时 $\Gamma(k)=\gamma^1(0)\gamma^1(1)\cdots\gamma^1(k-1)$，$k<0$ 时 $\Gamma(k)=\gamma^1(k)\gamma^1(k+1)\cdots\gamma^1(-1)$。对\sr{fig:13.3}{图~13.3}(ii) 作类似论证，得到
\[
\begin{aligned}
&\gamma^d(k_d,\ldots,k_{n+1},k_n,\ldots,k_1)\gamma^1(k_n)\\
&\qquad=\gamma^{d-1}(k_d,\ldots,k_n+k_{n+1},\ldots,k_1)\gamma^2(k_{n+1},k_n),
\end{aligned}
\]
{\looseness=-1
从而 $\gamma^d(0,\allowbreak \ldots,\allowbreak 0)=\gamma^2(0,\allowbreak 0)^{d-1}\Gamma(2-d)$。因此，$\F(\symbfit{L}^\#)^\to$ 与 $\widetilde{\F}(\symbfit{L}^\#)^\to$ 之间有一个简单的线性同构，它将每个 Floer 上链群 $CF^k$ 乘以 $\Gamma(k)\gamma^2(0,\allowbreak 0)$。这已足够满足我们的目的，因此不尝试显式确定剩下的符号。\par}
